\documentclass[11pt,a4paper]{article}
\usepackage[utf8]{inputenc}\usepackage[T1]{fontenc}\usepackage[italian,english]{babel}
\usepackage{amsmath,amssymb,amsthm}\usepackage[margin=2.2cm]{geometry}\usepackage{longtable,booktabs,array,calc}
\usepackage{listings,xcolor}\usepackage{hyperref}\usepackage{url}
\providecommand{\tightlist}{\setlength{\itemsep}{0pt}\setlength{\parskip}{0pt}}
\lstset{basicstyle=\ttfamily\scriptsize,breaklines=true,frame=single,captionpos=b,columns=fullflexible,keepspaces=true,inputencoding=utf8,extendedchars=true}
\title{Proof Pursuit --- Problem 2, part 5\\ \large prove, decisioni degli agenti, codice verificato, letteratura}
\author{Team Proof Pursuit (Researcher: T. Tumini; Referee: gabundos; Creative: M. Renzi; gio3r3gio)}
\date{\today}
\begin{document}\maketitle\tableofcontents
\section*{Come leggere questo rapporto}
Per ogni cella: la richiesta ufficiale, lo stato dichiarato dal team (mai ``risolto'' senza revisione umana), la consegna
con la prova, poi la traccia delle decisioni degli agenti (Researcher: famiglia e motivo; Referee: verdetto ed errore
fatale; Creative: quando e perché è intervenuto) e il codice su cui il tentativo poggia con l'esito della riesecuzione
fidata. DIMOSTRATO, CITATO, VERIFICATO SU INTERVALLO FINITO e CONGETTURA restano distinti. L'architettura del sistema
è descritta nel README del repository.
\section{Problem 2 — Uphill paths on the hypercube}
\subsection{Cella 5 (8 punti) — stato: parziale}
\subparagraph{\texorpdfstring{Parte 5 (C5) --- Bounds for
\(U(Q_9)\)}{Parte 5 (C5) --- Bounds for U(Q\_9)}}\label{parte-5-c5-bounds-for-uq_9}

\textbf{Punteggio:} 8 points · \textbf{Valutazione:} Judged

\(Q_9\) has \(512\) vertices and \(2304\) edges. The best bounds known
to the organisers are \[
2368 \le U(Q_9) \le 2400;
\] the lower bound is unpublished. Improve either one: prove that
\(U(Q_9) \ge 2369\), or exhibit a labelling of \(Q_9\) with at most
\(2399\) uphill paths.

\textbf{Enunciato protetto dato agli agenti (state.json).} \texttt{Improve one of the known bounds 2368 \textless{}= U(Q\_9) \textless{}= 2400: either prove U(Q\_9) \textgreater{}= 2369, or exhibit a labelling of Q\_9 with at most 2399 uphill paths (its paths are counted mechanically).}

\textbf{Evidenza registrata in STATUS.md.} la costruzione a stelle con un codice (8,20,3) dà esattamente 2400: il bound noto è questa famiglia; per scendere serve una foresta non a stelle; bozza parziale in inglese in submission/ (parziali.py)

\subsubsection{Consegna (prova)}
\paragraph{Problema 2 --- Parte 5 --- bozza di consegna
PARZIALE}\label{problema-2-parte-5-bozza-di-consegna-parziale}

\textbf{Stato dichiarato: PARZIALE.} Nessuna soluzione completa: qui
sotto ciò che è stabilito, la formalizzazione, la posizione rispetto
alla letteratura e ciò che resta aperto. Nulla è dichiarato dimostrato
oltre quanto scritto.

\subparagraph{1. Risultato e ambito}\label{risultato-e-ambito}

\textbf{Richiesta ufficiale.} \#\# Parte 5 (C5) --- Bounds for
\(U(Q_9)\)

\textbf{Punteggio:} 8 points · \textbf{Valutazione:} Judged

\(Q_9\) has \(512\) vertices and \(2304\) edges. The best bounds known
to the organisers are \[
2368 \le U(Q_9) \le 2400;
\] the lower bound is unpublished. Improve either one: prove that
\(U(Q_9) \ge 2369\), or exhibit a labelling of \(Q_9\) with at most
\(2399\) uphill paths.

\textbf{Cosa consegniamo.} Nessun miglioramento dei bound noti.
Contributo: la costruzione a stelle con un codice pari a distanza 4 di
taglia 20 dà esattamente il bound superiore noto; dimostrazione che
dentro questa famiglia non si può scendere.

\subparagraph{2. Dimostrazione}\label{dimostrazione}

Costruzione a stelle (vedi parte 3): totale \(=(d+1)2^{d-1}-(d-1)|R|\)
con \(R\) codice pari a distanza 4. Per \(d=9\), \(|R|\le A(8,3)=20\) e
il totale minimo della famiglia è \(2560-160=2400\), cioè il bound
superiore degli organizzatori. Quindi il bound noto è (con ogni
probabilità) proprio questa costruzione, e migliorarlo richiede una
foresta di non-picchi non a stelle (alberi più grandi, che riducano il
numero di componenti sotto 96). Impostato: cercare insiemi indipendenti
\(P\) di \(Q_9\) con \(|P|<236\) e complemento aciclico; non eseguito.

\subparagraph{3. Verifica: istruzioni, dipendenze,
tempi}\label{verifica-istruzioni-dipendenze-tempi}

Codice disponibile (Python 3, libreria standard; ogni script gira in
meno di un minuto): - \texttt{problema-2/certificati/costruzione\_88.py}
- \texttt{problema-2/certificati/q3\_esaustivo.py} -
\texttt{problema-2/certificati/q3\_verifica\_indipendente.py} -
\texttt{problema-2/certificati/q4\_branch\_and\_bound.py} -
\texttt{problema-2/certificati/q4\_ricerca\_locale.py} -
\texttt{problema-2/certificati/q4\_verifica\_etichettatura\_indipendente.py}
- \texttt{problema-2/certificati/q5\_stella\_verifica\_indipendente.py}
- \texttt{problema-2/certificati/verifica\_dfs\_88.py} -
\texttt{problema-2/certificati/verifica\_etichettatura\_q3.py} -
\texttt{problema-2/certificati/verifica\_etichettatura\_q4.py} -
\texttt{problema-2/certificati/verifica\_q5\_bipartita.py} -
\texttt{problema-2/certificati/verifica\_q5\_indipendenti.py} -
\texttt{problema-2/esperimenti/conta\_cammini.py} -
\texttt{problema-2/esperimenti/costruzione\_stelle.py} -
\texttt{problema-2/esperimenti/ricottura\_controllo.py}

\subparagraph{4. Fonti e contributo}\label{fonti-e-contributo}

Letteratura arXiv (ricerca deterministica
\texttt{tools/cerca\_letteratura.sh}, abstract letti, non usata come
prova): - arXiv:1412.3893v1 --- The competition between simple and
complex evolutionary trajectories in asexual populations (Ian E. Ochs,
Michael M. Desai, 2014); solo abstract letto. Nessun lavoro arXiv sui
cammini in salita sull'ipercubo: il problema risulta inedito.

\subparagraph{5. Limiti e parti
irrisolte}\label{limiti-e-parti-irrisolte}

Sia il bound inferiore sia quello superiore.

\subsubsection{Decisioni degli agenti}
\paragraph{Run \texttt{p2\_c5}}

\subsubsection{Codice su cui poggia il tentativo}
Nessun codice nei tentativi.

\section{arXiv literature consulted}
\begin{thebibliography}{99}
\bibitem{1412.3893v1} Ian E. Ochs, Michael M. Desai. \emph{The competition between simple and complex evolutionary trajectories in asexual populations}. arXiv:1412.3893v1 (2014). \url{http://arxiv.org/abs/1412.3893v1}
\end{thebibliography}
\end{document}
